% Local editorial examples following CP-VI-0078; canonical problem and solution remain unchanged.
\par\medskip
\noindent\textbf{Why the residue field is essential.}
If $s$ corresponds to $\mathfrak p\subset A$, the fiber is
\[
X_s=\operatorname{Spec}\!\bigl(B\otimes_A\kappa(\mathfrak p)\bigr),
\qquad
\kappa(\mathfrak p)=\operatorname{Frac}(A/\mathfrak p).
\]
In general $\kappa(\mathfrak p)$ is not $A/\mathfrak p$:
for the generic point of $\operatorname{Spec}k[t]$ it is $k(t)$.
The geometric fiber over $s$ is obtained by one further extension
of scalars to an algebraic closure $\overline{\kappa(\mathfrak p)}$.

\par\medskip
\noindent\textbf{Worked example 1 (a double point at the origin).}
Consider $f:\operatorname{Spec}k[x]\to\operatorname{Spec}k[t]$
induced by $t\mapsto x^2$. At $s=(t)$,
\[
X_s\cong\operatorname{Spec}\bigl(k[x]/(x^2)\bigr).
\]
Its only point has length two: the scheme-theoretic fiber records
multiplicity, whereas the set-theoretic inverse image only records
$x=0$.

\par\medskip
\noindent\textbf{Worked example 2 (a general and a generic fiber).}
For $s=(t-a)$ with $a\in k^\times$,
\[
X_s\cong\operatorname{Spec}\bigl(k[x]/(x^2-a)\bigr).
\]
When $\operatorname{char}k\ne2$ and $a$ is a nonzero square in $k$,
this is a disjoint pair of reduced $k$-points.  Over the generic
point $(0)$, in characteristic different from two, the coordinate
ring is $k(t)[x]/(x^2-t)$, a quadratic field extension of $k(t)$.

\par\medskip
\noindent\textbf{Worked example 3 (fiber dimension can jump).}
Take $A=k[t]$ and $B=k[t,x]/(tx)$, with its evident $A$-algebra
structure.  For $s=(t-a)$, $a\ne0$, one obtains
$B\otimes_A k\cong k[x]/(ax)\cong k$, a zero-dimensional fiber.
But at $s=(t)$,
\[
B\otimes_A k\cong k[x],
\]
so the special fiber is an affine line of dimension one.
A morphism need not have constant fiber dimension.

\par\medskip
\noindent\textbf{Worked example 4 (arithmetic splitting changes with $p$).}
Let $A=\mathbb Z$ and $B=\mathbb Z[x]/(x^2-2)$. At a prime $(p)$,
\[
B\otimes_A\mathbb F_p\cong
\mathbb F_p[x]/(x^2-\overline2).
\]
For $p=2$ the fiber is the non-reduced double point
$\operatorname{Spec}\mathbb F_2[x]/(x^2)$.
For $p=3$ the polynomial $x^2+1$ is irreducible and the fiber is
$\operatorname{Spec}\mathbb F_9$.
For $p=7$, $x^2-2=(x-3)(x+3)$ has distinct roots, so the
fiber is the union of two reduced $\mathbb F_7$-points.
These three fibers have different scheme structures even though
they arise from the same morphism.
